\scanpage{158}\mspage{13}
faisceau de coeff. sur $X$ annulé par $n$, la dualité globale nous dit que $Rf_{\deleted{\unsure{*}}!}(F)$ et $Rf_*(R\underline{\mathrm{Hom}}\,(F,\mu_n^{\otimes d}))$ sont duals l'un de l'autre à valeurs dans $(\mathbb{Z}/n\mathbb{Z})_k$, p.ex. $Rf_!(\mathbb{Z}/n\mathbb{Z})$ et $Rf_*(\mu_n^{\otimes d})$, ou $Rf_!(\mu_n^{\otimes d})$ et $Rf_*(\mathbb{Z}/n\mathbb{Z})$, sont duals l'un de l'autre -- du moins modulo un shift d'amplitude $2d$ sur les dimensions. (Comme $\mathbb{Z}/n\mathbb{Z}$ est injectif sur lui-même, cela donne d'ailleurs des dualités parfaites
\[
R^{i}f_{\deleted{*}!}(F)\times R^{2d-i}f_*(\mathbb{R}\underline{\mathrm{Hom}}\,(F,\mu_n^{\otimes d}))\longrightarrow\mathbb{Z}/n\mathbb{Z}\big{)}.
\]

Si maintenant $n$ n'est plus supposé premier à la caractéristique, \inserted{p.ex \inserted{est} une puissance de $p=\mathrm{car}\,k>0$,} il semble que tout \unsure{s'écroule|s'ecroule}, à commencer qu'on ne sait plus \inserted{(pour $d>1$)} par quoi remplacer $\mu_n^{\otimes d}$... Le miracle assez extraordinaire, c'est que pour $d=1$ i.e. $X$ une courbe lisse, tout continue à marcher parfaitement, pourvu qu'on l'énonce avec soin ! Les premières vérifications se font p.ex. avec $F=\mathbb{Z}/p\mathbb{Z}$, $\mu_p$, $\alpha_p$, $X$ complète -- on trouve que c'est OK grâce essentiellement à l'autodualité de la jacobienne. On peut sophistiquer les exemples, en prenant des coefficients $F$ \uline{tordus}, et $X$ non complète -- on se convainc que ça marche toujours.! Simplement il faut faire attention qu'ici les ${\unsure{R|\mathbb{R}}}^if_*(F)$, $R^if_!(F')$ ont une structure ``continue'' (ce sont essentiellement des \unsure{groupes} \unsure{ind-pro-algébriques)}
% archive sheet number: 158
